% id: 37-37
% book: 쎈B 미적분1 · 03 미분계수와 도함수
% section: 유형 09 구간에 따라 다르게 정의된 함수의 미분가능성
% figure: none
% answer: 
% answer_source: 
% variant_level: 0 (원본 전사)
% 이 파일은 build.mjs 가 items.json 에서 만든다. 고칠 때는 items.json 을 고친다.
\smash{\raisebox{6.5mm}{\dmkichul{쎈B 미적분1 37쪽 37번}}}
함수 $f(x)=\begin{cases} 7-2x & (x>3) \\ x+a & (x\le 3) \end{cases}$에 대하여 함수 $g(x)$를\par\smallskip\dispeq{$g(x)=\begin{cases} (x^2-3x)f(x) & (x>3) \\ (ax+b)f(x) & (x\le 3) \end{cases}$}\par\smallskip\noindent 로 정의할 때, 함수 $g(x)$는 $x=3$에서 미분가능하다. 정수 $a$, $b$에 대하여 $a+b$의 값을 구하시오.
